\documentclass[11pt]{article}
\usepackage[margin=1in]{geometry}
\usepackage{booktabs}
\usepackage{longtable}
\usepackage{hyperref}
\usepackage{amsmath}
\title{WEC2026 Goal-Scoring Prediction: Data, Preprocessing, and Modeling Report}
\author{Automated experiment report}
\date{\today}

\begin{document}
\maketitle

\section{Executive Summary}
The task is a rare-event binary classification problem: predicting whether a player scores after a match checkpoint.
The modeling table contains 3,486 checkpoint observations, 203 positives (5.82%), 31 fixtures, and 869 player appearances.
The best grouped cross-validated model in this run was \textbf{catboost} on \textbf{all\_features} features with ROC-AUC 0.676, PR-AUC 0.124, and balanced accuracy 0.640.
AutoGluon was not available in this run.

\section{Data Audit}
The checkpoint file provides player state every 15 minutes, including position, home status, formation, substitution times, physical load, shots, and the target \texttt{scored\_after}.
Four event files were joined using \texttt{player\_appearance\_id} and absolute match time: passes, high-speed runs/sprints, shots, and actions under defensive pressure.
The target is highly imbalanced; therefore ROC-AUC, PR-AUC, balanced accuracy, Brier score, and log loss are more informative than raw accuracy.

\section{Preprocessing}
All event-derived features were computed with temporal filters:
\[
  \texttt{minute\_in} \leq \texttt{event\_abs\_minute} \leq \texttt{checkpoint\_abs\_minute}.
\]
The implementation converts period-relative minutes to absolute match minutes so second-half and extra-time checkpoints are handled correctly.
Reference quantities such as expected-threat zone values, pressure turnover baselines, and positional baselines were fitted inside each cross-validation training fold before being applied to validation checkpoints.
Categorical variables were one-hot encoded from training folds only, and numerical missing values were median-imputed from training folds only.

\section{Feature Families}
The all-feature model uses signals from every available file:
physical intensity and shots from \texttt{players\_quarters\_final.csv}, run event intensity from \texttt{player\_appearance\_run.csv}, pass volume and xT proxies from \texttt{player\_appearance\_pass.csv}, shot context from \texttt{player\_appearance\_shot\_limited.csv}, and pressure resistance plus pressing-applied metrics from \texttt{player\_appearance\_behaviour\_under\_pressure.csv}.
Recent-vs-cumulative ratios were included to test short-term intensity surges.

\section{Model Selection and HPO}
The benchmark used fixture-grouped stratified cross-validation to prevent match-level leakage.
Models attempted in this practical run included logistic regression, ExtraTrees, RandomForest, XGBoost, CatBoost, LightGBM, TabPFN, stacking, and optional AutoGluon.
TabPFN probabilities were included as a base learner in the stacked ensemble whenever the package was available, rather than being evaluated only as a standalone benchmark.
Skipped or unavailable models: tabpfn: TabPFN requires a one-time license acceptance to download.
XGBoost used Optuna HPO when the requested trial budget was nonzero; other model parameters were conservative regularized defaults appropriate for the small sample size.

\section{Grouped Cross-Validated Results}
\begin{longtable}{llrrrrr}
\toprule
Model & Feature set & ROC-AUC & PR-AUC & Bal. Acc. & F1 & Brier \\
\midrule
catboost & all\_features & 0.676 & 0.124 & 0.640 & 0.206 & 0.181 \\
stacked & all\_features & 0.671 & 0.121 & 0.630 & 0.177 & 0.054 \\
xgboost & all\_features & 0.692 & 0.120 & 0.500 & 0.000 & 0.061 \\
extra\_trees & all\_features & 0.668 & 0.117 & 0.500 & 0.000 & 0.061 \\
xgboost & base\_checkpoint & 0.677 & 0.112 & 0.500 & 0.000 & 0.062 \\
xgboost & base\_plus\_passing & 0.674 & 0.112 & 0.500 & 0.000 & 0.062 \\
xgboost & cumulative\_only & 0.643 & 0.111 & 0.506 & 0.028 & 0.069 \\
xgboost & base\_plus\_pressure & 0.662 & 0.107 & 0.500 & 0.000 & 0.062 \\
logistic & all\_features & 0.619 & 0.100 & 0.589 & 0.150 & 0.197 \\
xgboost & recent\_only & 0.637 & 0.092 & 0.500 & 0.000 & 0.068 \\
xgboost & sprints\_shots & 0.618 & 0.085 & 0.504 & 0.026 & 0.081 \\
xgboost & external\_context & 0.552 & 0.074 & 0.547 & 0.131 & 0.167 \\
\bottomrule
\end{longtable}

\section{AutoGluon Holdout Result}
\begin{longtable}{llrrrrr}
\toprule
Model & Feature set & ROC-AUC & PR-AUC & Bal. Acc. & F1 & Brier \\
\midrule
No AutoGluon holdout result was recorded. \\
\bottomrule
\end{longtable}

\section{Ablation Results}
\begin{longtable}{lrrrr}
\toprule
Feature set & ROC-AUC & PR-AUC & Bal. Acc. & F1 \\
\midrule
all\_features & 0.692 & 0.120 & 0.500 & 0.000 \\
base\_checkpoint & 0.677 & 0.112 & 0.500 & 0.000 \\
base\_plus\_passing & 0.674 & 0.112 & 0.500 & 0.000 \\
cumulative\_only & 0.643 & 0.111 & 0.506 & 0.028 \\
base\_plus\_pressure & 0.662 & 0.107 & 0.500 & 0.000 \\
recent\_only & 0.637 & 0.092 & 0.500 & 0.000 \\
sprints\_shots & 0.618 & 0.085 & 0.504 & 0.026 \\
external\_context & 0.552 & 0.074 & 0.547 & 0.131 \\
\bottomrule
\end{longtable}

\section{Feature Importance}
\begin{longtable}{lr}
\toprule
Feature & PR-AUC drop after permutation \\
\midrule
cumul\_pass\_bottom\_share & 0.016 \\
cumul\_pressure\_quality & 0.006 \\
is\_defender & 0.004 \\
cumul\_pressure\_top\_share & 0.004 \\
last15\_pass\_bottom\_share & 0.004 \\
cumul\_pass\_top\_share & 0.003 \\
cumul\_pass\_count & 0.002 \\
cumul\_pass\_receiver\_diversity & 0.001 \\
last15\_pass\_top\_share & 0.001 \\
checkpoint\_abs\_min & 0.001 \\
cumul\_positive\_xt\_count & 0.001 \\
recent\_hsr\_ratio & 0.001 \\
\bottomrule
\end{longtable}

\section{Research Questions}
\paragraph{RQ1.}
Player behaviour during a match is predictive but noisy because goals are rare. In this run the best grouped-CV model reached ROC-AUC 0.676, PR-AUC 0.124, and balanced accuracy 0.640.
\paragraph{RQ2.}
The strongest determinants are identified by a combination of ablation deltas and permutation importance. The leading stable signals in this run were: cumul\_pass\_bottom\_share, cumul\_pressure\_quality, is\_defender, cumul\_pressure\_top\_share, last15\_pass\_bottom\_share, cumul\_pass\_top\_share.
\paragraph{RQ3.}
The sprint-and-shot-only ablation should be compared against the all-feature model. Its result in this run was 0.085 PR-AUC, so it is not assumed sufficient unless it matches the all-feature score.
\paragraph{RQ4.}
Passing and pressure data are useful when their ablation improves over the base checkpoint model. The report tables show whether \texttt{base\_plus\_passing} and \texttt{base\_plus\_pressure} improve ROC-AUC/PR-AUC/balanced accuracy.
\paragraph{RQ5.}
The \texttt{recent\_only} and \texttt{cumulative\_only} ablations compare short-term and accumulated performance. The stronger family is the one with higher grouped-CV PR-AUC and ROC-AUC.
\paragraph{RQ6.}
Recent-vs-cumulative ratios and intensity-surge interactions directly test whether short-term intensity above a player's own baseline increases scoring probability.
\paragraph{RQ7.}
External factors were modeled through home status, period, checkpoint time, extra time, formation, position, and substitution context. Their influence is assessed by the \texttt{external\_context} ablation and permutation importance.

\section{Limitations}
Only 203 positive examples are available, so fold-level variance is expected. Some event locations are coarse three-zone labels, which limits xT precision. The stacked model is evaluated with grouped OOF meta-predictions, but it should still be validated on a hidden competition set before final claims.

\end{document}
